% Guarding Walls Contiguously -- AI journal
% INFO-F420 Computational Geometry, Universite libre de Bruxelles, 2026-2027
% Compile with pdfLaTeX (Overleaf: default settings). No image needed.

\documentclass[11pt,a4paper]{article}

\usepackage[T1]{fontenc}
\usepackage[utf8]{inputenc}
\usepackage{lmodern}
\usepackage{microtype}
\usepackage[a4paper,margin=2.2cm]{geometry}
\usepackage{amsmath,amssymb}
\usepackage{booktabs}
\usepackage{longtable}
\usepackage{array}
\usepackage[shortlabels]{enumitem}
\usepackage{xcolor}
\usepackage[most]{tcolorbox}
\usepackage[colorlinks=true,linkcolor=blue!50!black,citecolor=blue!50!black,urlcolor=blue!50!black]{hyperref}

\setlist{itemsep=3pt,topsep=4pt}
\renewcommand{\arraystretch}{1.35}
\newcolumntype{L}[1]{>{\raggedright\arraybackslash}p{#1}}

\title{\textbf{AI Journal}\\[4pt]
\large Guarding Walls Contiguously\\[2pt]
\normalsize INFO-F420 Computational Geometry, Universit\'e libre de Bruxelles, 2026--2027}
\author{Ilias Akajou (520466) \and Jomaa Abdulrahman (583629)}
\date{November 2026}

\begin{document}
\maketitle

%======================================================================
\section{Tools, Environment, and Workflow}
%======================================================================

We used \textbf{Claude Code} (Anthropic, model \textbf{Claude Opus 5.5}) in its desktop application, between 7 and 24 November 2026, in work sessions
led by Ilias.

\begin{itemize}[leftmargin=*]
\item \textbf{Interactive workspace:} the assistant had read/write access to the project folder. It could write and run the JavaScript code and the
automated tests, and open the website in a browser.
\item \textbf{How a session worked:} we stated what we wanted (a feature, an explanation, a check), the assistant produced it, and we reviewed the result in
the browser or in the generated documents before asking for the next change.
\item \textbf{Our rule:} every result we kept was checked, either by the automated tests or against the papers~\cite{deBerg2025,deBerg2026,Biniaz2025}, and each of
us is responsible for understanding and defending his half of the proof at the oral (Section~\ref{sec:ours}).
\end{itemize}

%======================================================================
\section{Our Work, Decisions, and Division of Labor}
\label{sec:ours}
%======================================================================

\begin{tcolorbox}[breakable,colback=green!4,colframe=green!45!black,boxrule=0.5pt,arc=2mm,left=2.5mm,right=2.5mm,top=1.5mm,bottom=1.5mm]
\begin{itemize}[leftmargin=*]
\item \textbf{Topic selection and project proposal.} Together, we selected the Contiguous Art Gallery Problem, chose the main paper~\cite{deBerg2025}, and wrote
our proposal (19 October), which structures the work into four parts with a bibliography. Following the professor's feedback, we revised the plan: Part~1
kept short, both of us on the core algorithm in Part~2, and one interactive widget for each key step of the proof.
\item \textbf{Architecture of the website.} The pedagogical structure is our decision. After reviewing a first version that felt too dense for a live
presentation, Ilias asked for a streamlined \emph{Overview} page containing only what we present orally, with the proof chapters, the interactive figures and
the game on separate pages.
\item \textbf{The simple-polygon assumption.} Reviewing the site, Jomaa noticed that it did not state that the polygon must be \emph{simple} (no holes, no
self-intersections), although the algorithm fails otherwise. We had this checked in the papers and explained on every page.
\item \textbf{Peer feedback from the course forum.} We posted a preliminary version on the forum and acted on a classmate's remark: the animations of the
overview now play once, can be replayed, and can be switched off.
\item \textbf{Division of the proof, the code and the oral.} Following the professor's feedback, we split the proof into two halves:
  \begin{itemize}[leftmargin=1.2em]
  \item \textbf{Ilias (first half, Sections~2.1--2.4, Figures~2--5):} visibility cores, the greedy algorithm and its $k+1$ bound, good and bad guards, and
  the replacement of guards by arrangement vertices; also the problem (overview and Figure~1) and the game.
  \item \textbf{Jomaa (second half, Sections~2.5--2.10, Figures~6--11):} the candidate guards $G_P$ and starting points $U_P$, Theorem~5.4, maximal visible
  chains (Lemma~6.2), one greedy step, the complete algorithm and the $\Theta(n\log n)$ algorithm~\cite{deBerg2026}; also the automated tests.
  \end{itemize}
\item \textbf{Publication and final checks.} Ilias created the public GitHub repository and published the site on GitHub Pages. We had the website and the
report checked against the papers, the professor's feedback and the UV course page (Section~\ref{sec:checks}).
\end{itemize}
\end{tcolorbox}

%======================================================================
\section{What the AI Did, at Our Request}
%======================================================================

\begin{itemize}[leftmargin=*]
\item \textbf{Understanding the papers:} the assistant read the papers~\cite{Biniaz2025,Merrild2024,Biniaz2024,Robson2024,deBerg2025,deBerg2026} and wrote
explanations for us in French: the subject from scratch, the whole project in detail with questions to test ourselves, and one presentation guide per
student with likely questions. We used them to learn the definitions and the logic of the proof, and to prepare the oral.
\item \textbf{Building the website:} following the structure we chose, the assistant designed and built the website and implemented, in JavaScript, the
algorithm of~\cite{deBerg2025}, the eleven interactive figures and the game. We reviewed each version and asked for changes (structure, animations,
simple-polygon notice, phone layout).
\item \textbf{Testing:} the assistant wrote automated tests that compare every step with slower brute-force computations, and used them to find and correct
bugs (Section~\ref{sec:checks}).
\item \textbf{Texts and compliance:} the assistant drafted the texts of the website and of the report, and its LaTeX version, and checked them against the
papers and the course page.
\end{itemize}

%======================================================================
\section{Significant Prompts and How We Used the Responses}
\label{sec:prompts}
%======================================================================

Our main requests (written in French, summarised here), the response of the assistant, and how we used it.

{\small
\begin{longtable}{@{}L{2.7cm} L{5.8cm} L{7.2cm}@{}}
\toprule
Objective & What we asked & Response \& how we used it \\
\midrule
\endhead

\textbf{Building}\newline 07/11/2026 &
Implement the project of our proposal, following the professor's feedback, and give us documents that explain the papers. &
A geometry library, the algorithm, the tests and a first website with 11 interactive figures and a game; explanations in French that we used to start
understanding the proof. \\

\textbf{Structure}\newline 09/11/2026 &
Restructure the site: an overview with only what we will present orally, the details on other pages, and one document that explains everything. &
The current structure (overview, five proof chapters, game, report) and a complete explanation document. The overview became the support of our
presentation. \\

\textbf{Understanding}\newline 11/11/2026 &
Explain the subject to us from scratch, in French: we have to understand it in order to explain it ourselves. &
A step-by-step explanation, from the basic definitions to the main theorem. We used it to learn the vocabulary (chains, visibility cores, good and bad
guards) and the logic of the proof. \\

\textbf{Verification}\newline 11/11/2026 &
Jomaa's remark: the site does not say that the polygon must be simple, and the algorithm fails otherwise. Check it, and check the course requirements. &
Confirmed in the papers: Merrild et al.~\cite{Merrild2024} state that their approach does not work with holes and that this case is open. The assumption
is now explained on every page, with the reasons why it is needed. \\

\textbf{Organisation}\newline 23/11/2026 &
Who explains which part of the proof at the oral, and how: one guide per student. &
A work-distribution section on the site and one presentation guide per student, with the text of each part and likely questions; we use them to
rehearse. \\

\textbf{Organisation}\newline 23/11/2026 &
State precisely who did what, following the professor's feedback. &
The team section of the report: for each step of the proof, its section, its figure, its code and who is responsible for it. \\

\textbf{Understanding}\newline 24/11/2026 &
Explain the whole project to us in detail, in French: the course requirements, every step of the proof, the code and the site. &
A detailed document with, for each step, a one-sentence summary, an example and what to remember, and questions to test ourselves. \\

\textbf{Verification}\newline 24/11/2026 &
Check that nothing is wrong: compare the report with the papers, the professor's feedback, the site and the course page on the UV. &
The corrections listed in Section~\ref{sec:checks}, a bibliography following the citation guidelines, and instructions to run the site from a local file. \\

\bottomrule
\end{longtable}
}

%======================================================================
\section{How We Checked, Corrected, and Rejected AI Outputs}
\label{sec:checks}
%======================================================================

\begin{itemize}[leftmargin=*]
\item \textbf{Automated brute-force verification (1,541 / 1,541 checks passing):} visibility is compared with direct segment tests, $\mathrm{next}(u)$ with a
dense grid of guard positions, and the optimum with the greedy algorithm started from 200 evenly spaced points (no run may use more than $k+1$ guards).
Every construction of the proof was also checked on the guards of 60 random polygons: 8,647 good guards and 1,381 bad guards were replaced as the lemmas
predict, Lemma~6.2 matched the direct computation in 5,014 cases, and Theorem~5.4 ended at a point of $U_P$ in 61 cases.
\item \textbf{The Lemma~6.2 bug:} the first implementation took the intersection point closest to the guard~$g$ and disagreed with the direct computation in
about 19\% of random cases. The proof of~\cite{deBerg2025} defines the point as the one closest to~$m$; with this rule, all tests pass. The result of the
paper is correct, but its statement needs this precision when a chain winds around $g$ more than once.
\item \textbf{Verification against the papers:} comparing the report with the papers corrected four errors: a running time wrongly attributed to Robson
et al.~\cite{Robson2024}, the description of Biniaz et al.~\cite{Biniaz2024}, the page count of~\cite{deBerg2025}, and the statement of Lemma~6.8. On the
website, it also corrected a wrong number on the overview and aligned the angle notation with the paper.
\item \textbf{Course checklist and citation guidelines:} following the UV page on citations, the report now cites the journal version of Hershberger and Suri
(1995) instead of the conference version (SODA 1993) and gives DOIs, which we checked; the submission explains how to run the site from a local file.
\item \textbf{Rejected:}
  \begin{itemize}[leftmargin=1.2em]
  \item a claim in a figure (``being in the core and seeing both ends implies seeing the whole chain''), removed because it could not be justified in all cases;
  \item an observation from our experiments (starting the greedy algorithm at the vertices was always enough in our tests), reported only as an observation,
  since nothing in the papers guarantees it;
  \item a reference to a workshop version of~\cite{deBerg2026}, not added because we could not verify its details.
  \end{itemize}
\end{itemize}

%======================================================================
\begin{thebibliography}{9}
%======================================================================

\bibitem{Biniaz2025}
A.~Biniaz, A.~Maheshwari, M.~C.~R. Merrild, J.~S.~B. Mitchell, S.~Odak, V.~Polishchuk, E.~W. Robson, C.~M. Rysgaard, J.~K.~R. Schou, T.~Shermer,
J.~Spalding-Jamieson, R.~Svenning, D.~W. Zheng.
\newblock Polynomial-time algorithms for contiguous art gallery and related problems.
\newblock In \emph{Proc.\ 41st International Symposium on Computational Geometry (SoCG 2025)}, volume 332 of LIPIcs, pages 20:1--20:21, 2025.
\newblock \href{https://doi.org/10.4230/LIPIcs.SoCG.2025.20}{doi:10.4230/LIPIcs.SoCG.2025.20}.

\bibitem{Merrild2024}
M.~C.~R. Merrild, C.~M. Rysgaard, J.~K.~R. Schou, R.~Svenning.
\newblock The contiguous art gallery problem is solvable in polynomial time.
\newblock \href{https://arxiv.org/abs/2412.13938}{arXiv:2412.13938}, 2024 (v3, 2025). Merged into~\cite{Biniaz2025}; full version.

\bibitem{Biniaz2024}
A.~Biniaz, P.~Bose, A.~Maheshwari, J.~S.~B. Mitchell, S.~Odak, V.~Polishchuk, T.~C. Shermer.
\newblock Contiguous boundary guarding.
\newblock \href{https://arxiv.org/abs/2412.15053}{arXiv:2412.15053}, 2024 (v2, 2025). Merged into~\cite{Biniaz2025}; full version.

\bibitem{Robson2024}
E.~W. Robson, J.~Spalding-Jamieson, D.~W. Zheng.
\newblock The analytic arc cover problem and its applications to contiguous art gallery, polygon separation, and shape carving.
\newblock \href{https://arxiv.org/abs/2412.15567}{arXiv:2412.15567}, 2024 (v2, 2025). Merged into~\cite{Biniaz2025}; full version.

\bibitem{deBerg2025}
S.~de~Berg, J.~Conradi, I.~van~der~Hoog, F.~Staals.
\newblock Simpler and faster contiguous art gallery.
\newblock \href{https://arxiv.org/abs/2508.09734}{arXiv:2508.09734}, 2025.

\bibitem{deBerg2026}
S.~de~Berg, J.~Conradi, I.~van~der~Hoog, E.~Rotenberg.
\newblock The contiguous art gallery problem is in $\Theta(n\log n)$.
\newblock \href{https://arxiv.org/abs/2511.02960}{arXiv:2511.02960}, 2025 (v6, 2026).

\end{thebibliography}

\end{document}
